\scanpage{36}
\mspage{\handwritten{6}}
2.\uline{Accouplements définis par une biextension}.
2.1. Soient $P$ et $Q$ deux \deleted{g}Groupes abéliens sur $T$, et $n\handwritten{>}0$ un entier. Rappelons qu'on désigne par $\mathord{}_nP$ l\deleted{a}e noyau de la multiplication par $n\,\mathrm{id}_P$. Nous allons définir un homomorphisme canonique
\[ \handwritten{\alpha}=\handwritten{\alpha}^{\deleted{\unsure{N|n}}n}_{P,Q}:\quad {}_nP\otimes{}_nQ\mathbin{\handwritten{\longrightarrow}}\underline{\mathrm{Tor}}_1(P,Q),\tag*{(2.1.1)} \]
qui sera le composé
\[ {}_nP\otimes{}_nQ\mathbin{\handwritten{\longrightarrow}}\underline{\mathrm{Tor}}_1({}_nP,{}_nQ)\mathbin{\handwritten{\longrightarrow}}\underline{\mathrm{Tor}}_1(P,Q),\tag*{(2.1.2)} \]
où la première flèche est $\handwritten{\alpha}^n_{{}_nP,{}_nQ}$, et la deuxième est celle déduite des inclusions $\mathord{}_nP\handwritten{\to}P$, $\mathord{}_nQ\handwritten{\to}Q$. On est donc ramené à définir (2.1.1) dans le cas où on a $nP=0,nQ=0$, i.e. où $P$ et $Q$ sont des $\underline{\mathrm Z}/n\underline{\mathrm Z}$ - Modules.
\deleted{Calculant alors $\underline{\mathrm{Tor}}_{\unsure{1|i}}(P,Q)$ à l'aide d'une} Supposons seulement \deleted{\unsure{$n$|$nP$}}\inserted{$nQ=0$.} Prenant une résolution plate $\deleted{\unsure{M}}L\mathord.$ de $\deleted{Q}P$
\deleted{\illegible{1}}
(en tant que $\underline{\mathrm Z}$-Module), \deleted{comme $\underline{\mathrm H}_{\unsure{1}}(P\otimes M\mathord.)$ \unsure{$=$}} on trouve des isomorphismes canoniques
\[ \deleted{\mathrm H\mathord.(L\mathord.\otimes Q)=\mathrm H\mathord.((L\mathord.\otimes_{\underline{\mathrm Z}}}\qquad\deleted{P\otimes} \]
\[ P\overset{\underline L}{\otimes}_{\underline{\mathrm Z}}Q\handwritten{\simeq}L\mathord.\otimes_{\underline{\mathrm Z}}Q\handwritten{\simeq}(L\mathord.\otimes_{\underline{\mathrm Z}}\underline{\mathrm Z}/n\underline{\mathrm Z})\otimes_{\underline{\mathrm Z}/n\underline{\mathrm Z}}Q, \]
qu'on peut écrire comme un isomorphisme dans la catégorie dérivée
\[ P\overset{\underline L}{\otimes}_{\underline{\mathrm Z}}Q\handwritten{\simeq}(P\overset{\underline L}{\otimes}_{\underline{\mathrm Z}}\underline{\mathrm Z}/n\underline{\mathrm Z})\overset{\underline L}{\otimes}_{\underline{\mathrm Z}/n\underline{\mathrm Z}}Q,\tag*{(2.1.3)} \]
donnant naissance à \deleted{\unsure{une}} \handwritten{la} suite spectrale de type homologique \inserted{\handwritten{“de Künneth”}} : \margin{\handwritten{(cf. SGA 1 $\overline{\underline{\mathrm{IV}}}$ 5.2)}}
\[ \underline{\mathrm{Tor}}^{\underline{\mathrm Z}}_*(P,Q)\handwritten{\Longleftarrow}E^2_{p,q}=\underline{\mathrm{Tor}}^{\underline{\mathrm Z}/n\underline{\mathrm Z}}_{\deleted{\unsure{q}}p}(\underline{\mathrm{Tor}}^{\underline{\mathrm Z}}_q(P,\underline{\mathrm Z}/n\underline{\mathrm Z}),Q).\tag*{(2.1.4)} \]
Rappelons d'ailleurs qu'on a des isomorphismes canoniques
\[ \underline{\mathrm{Tor}}^{\underline{\mathrm Z}}_0(P,\underline{\mathrm Z}/n\underline{\mathrm Z})=P_n\mathrel{\overset{\mathrm{dfn}}{=}}P/nP,\quad \underline{\mathrm{Tor}}^{\underline{\mathrm Z}}_1(P,\underline{\mathrm Z}/n\underline{\mathrm Z})\simeq {}_nP,\quad \underline{\mathrm{Tor}}^{\underline{\mathrm Z}}_i(P,\underline{\mathrm Z}/n\underline{\mathrm Z})=0\text{ si }i\geq2, \]
déduits du calcul des $\underline{\mathrm{Tor}}^{\underline{\mathrm Z}}_i(P,\underline{\mathrm Z}/n\underline{\mathrm Z})$ à l'aide de la résolution plate canonique suivante de $(\underline{\mathrm Z}/n\underline{\mathrm Z})_T$ :
\[ \begin{aligned}0\deleted{\unsure{n}}&\longrightarrow\underline{\mathrm Z}_T\overset{n}{\longrightarrow}\underline{\mathrm Z}_T\longrightarrow 0\deleted{\unsure{1}}\\ &\hspace{2.6em}\overset{\prime\prime}{\underline L_1}\qquad\overset{\prime\prime}{\underline L_0}\qquad,\unsure{\prime\prime}\end{aligned} \]
l'augmentation $\underline L_0\to(\underline{\mathrm Z}/n\underline{\mathrm Z})_T$ étant définie en envoyant 1 sur 1. Par suite, la suite spectrale (2.1.4) se réduit essentiellement à une suite exacte
\[ 0\to\underline{\mathrm{Tor}}^{\underline{\mathrm Z}/n\underline{\mathrm Z}}_2(P_n,\deleted{\unsure{*}}Q)\to{}_nP\otimes Q\mathbin{\overset{\handwritten{\alpha}}{\handwritten{\longrightarrow}}}\underline{\mathrm{Tor}}^{\underline{\mathrm Z}}_1(P,Q)\to\underline{\mathrm{Tor}}^{\underline{\mathrm Z}/n\underline{\mathrm Z}}_1(P_n,Q)\to0,\tag*{(2.1.5)} \]
où on a posé $P_n=P\otimes_{\underline{\mathrm Z}}\underline{\mathrm Z}/n\underline{\mathrm Z}=P/nP$. L'homomorphisme $\handwritten{\alpha}$ \deleted{\unsure{est}} cherché est \deleted{celui} \inserted{\handwritten{l'homomorphisme-coin $E^2_{0,1}\to\underline{\mathrm{Tor}}^{\mathbb Z}_1$}} qui figure dans la dernière suite exacte. Cette dernière étant évidemment fonctorielle en $P$ et en $Q$, la fonctorialité en $P$ montre que l'homomorphisme envisagé est bien le composé
\[ \handwritten{{}_n}P\otimes Q\to\underline{\mathrm{Tor}}_1\handwritten{({}_nP,Q)\longrightarrow\underline{\mathrm{Tor}}_1(P,Q),}\qquad\deleted{\unsure{N|n}} \]
% archive sheet number: 35
